\chapter{Przegląd metod wykrywania niechcianej poczty}

Praktyczny filtr pocztowy korzysta nie tylko z tematu i treści wiadomości, ale
też z nagłówków SMTP, adresów IP, domen, historii nadawcy, wyników
uwierzytelniania i informacji o odsyłaczach. Analiza tekstu jest jedną z jego
warstw. Pozostałe mechanizmy oceniają nadawcę, wykrywają znane kampanie albo
sprawdzają podejrzane adresy URL. W przeglądzie uwzględniono zarówno te
rozwiązania, jak i metody możliwe do uruchomienia na zbiorach tekstowych użytych
w części badawczej.

\section{Wielowarstwowe systemy filtrowania poczty}

W produkcyjnych skrzynkach pocztowych pierwsza decyzja często zapada jeszcze przed
analizą pełnej treści wiadomości. Dostawca poczty może sprawdzić, czy nadawca
spełnia wymagania techniczne, czy wiadomość przechodzi uwierzytelnianie SPF,
DKIM i DMARC, czy adres IP ma dobrą reputację oraz czy domena nie występuje w
kampaniach nadużyć. Wymagania publikowane dla nadawców poczty Gmail wskazują
wprost na znaczenie SPF lub DKIM, polityki DMARC, poprawnego formatu wiadomości
i utrzymywania niskiego odsetka zgłoszeń spamu~\autocite{gmail-sender-guidelines}.
Podobny obraz widać w ustawieniach Microsoft Defender for Office 365, gdzie
osobno konfigurowane są między innymi polityki antyspamowe, antyphishingowe,
obsługa podszywania się pod nadawcę i reakcje na DMARC~\autocite{microsoft-defender-o365-security-settings}.

SpamAssassin opiera decyzję na zestawie testów punktowanych, regułach
treści, listach dozwolonych i blokowanych, testach sieciowych oraz mechanizmach
uczących się~\autocite{spamassassin-conf}. Rspamd udostępnia osobne moduły dla
SPF, DKIM, DMARC, reputacji, list DNS, fuzzy hashy, phishingu, sieci neuronowych
i wielu innych sygnałów~\autocite{rspamd-modules}. Oba filtry działają jako
systemy punktowe, które składają wiele sygnałów w jedną decyzję.

\section{Metody regułowe, reputacyjne i uwierzytelnianie nadawcy}

W filtrze pocztowym reguła jest jawnie zapisanym warunkiem sprawdzanym na
wiadomości lub jej metadanych. Po spełnieniu warunku filtr może dodać określoną
liczbę punktów do wyniku spamu, zablokować wiadomość albo przekazać ją do dalszej
analizy. Reguły są nadal używane, ponieważ działają szybko, są interpretowalne i
pozwalają reagować na znane schematy nadużyć bez ponownego trenowania modelu.
Reguła może dotyczyć fragmentu nagłówka, nietypowego adresu URL, słowa często
występującego w danej kampanii albo kombinacji kilku prostszych testów.
Administrator może też podnieść lub obniżyć wagę konkretnego testu, co ułatwia
bezpośrednią kontrolę nad polityką filtrowania.

Słabą stroną reguł jest zależność od wcześniej zauważonych wzorców. Nadawcy
spamu mogą zmieniać pisownię słów, przenosić treść do obrazów, rotować domeny
lub modyfikować szablony wiadomości. Dlatego reguły zwykle nie występują same,
lecz są łączone z reputacją nadawcy, listami domen, uwierzytelnianiem i
klasyfikatorami treści. Przeglądy metod antyphishingowych również pokazują, że
obrona przed phishingiem obejmuje kilka rodzin podejść: listy znanych adresów,
heurystyki, analizę treści, analizę wizualną oraz metody uczenia maszynowego
~\autocite{wood-antiphishing-review-2022}.

Uwierzytelnianie nadawcy rozwiązuje inny problem niż klasyfikacja tekstu.
Mechanizmy SPF, DKIM i DMARC pomagają ustalić, czy wiadomość mogła zostać
wysłana w imieniu deklarowanej domeny. Poprawnie uwierzytelniona wiadomość może
być niechcianą reklamą albo pochodzić z przejętego konta, dlatego analiza treści
pozostaje potrzebna.

\section{Klasyczne metody uczenia maszynowego}

Klasyczne modele uczenia maszynowego oparte na cechach leksykalnych są
najprostszym porównaniem dla klasyfikacji treści. Wiadomość zamienia się wtedy
na wektor cech, na przykład przez zliczanie słów, n-gramy albo TF-IDF, a
następnie trenuje klasyfikator taki jak Naive Bayes, regresja logistyczna lub
SVM. Jest to stosunkowo proste podejście, które zapewnia tani trening i szybką
inferencję, a przy odpowiednio dobranych danych może osiągać dobre wyniki.

W pracy \emph{Learning to Filter Spam E-Mail: A Comparison of a Naive Bayesian
and a Memory-Based Approach} porównano Naive Bayes z podejściem pamięciowym.
Automatycznie uczone filtry wyraźnie przewyższały ręcznie konstruowane reguły
słów kluczowych~\autocite{androutsopoulos-naive-bayes-spam}. Praca
\emph{Classification of Spam Emails through Hierarchical Clustering and
Supervised Learning} obejmowała porównanie reprezentacji TF-IDF i bag-of-words
z klasyfikatorami Naive Bayes, drzewami decyzyjnymi oraz SVM
~\autocite{janez-martino-spam-supervised-2020}.

Porównanie z klasycznymi metodami pozwala ocenić, ile z wyniku można uzyskać
dzięki dopasowaniu do słownictwa zbioru treningowego. Wynik TF-IDF bliski
wynikowi modelu językowego podważałby zasadność użycia droższego rozwiązania.
Przewaga modelu językowego na danych z innych źródeł wskazywałaby natomiast na
lepszą generalizację.

Reprezentacja leksykalna dobrze wychwytuje powtarzalne zwroty, nazwy produktów,
typowe frazy oszustw i często występujące domeny. Gorzej radzi sobie z
parafrazą, długim kontekstem i wiadomościami, w których znaczenie wynika z
relacji między kilkoma zdaniami, ponieważ warianty typu bag-of-words tracą
kolejność słów i słabiej opisują zależności semantyczne
~\autocite{le-distributed-representations-2014}. Rzadkie cechy oraz krótka treść
wiadomości mogą dodatkowo prowadzić do przeuczenia i słabszej generalizacji
~\autocite{xu-alternative-text-representation-2013}. W danych e-mailowych oznacza
to ryzyko dopasowania się do cech konkretnego korpusu, na przykład
charakterystycznego formatowania albo powtarzalnych fragmentów stopki.

\section{Szybkie modele tekstowe i reprezentacje semantyczne}

Pomiędzy klasycznym TF-IDF a pełnym dostrajaniem transformera znajdują się metody
pośrednie. Jedną z nich jest \texttt{fastText}, który uczy reprezentacje słów z
n-gramów znakowych i stosuje liniowy klasyfikator. Reprezentacje słów
występujących w dokumencie są łączone przed wykonaniem klasyfikacji. Autorzy
\texttt{fastText} pokazali, że taki model może być konkurencyjny wobec
głębszych klasyfikatorów tekstu, a jednocześnie trenuje się i działa bardzo
szybko~\autocite{joulin-fasttext-bag-of-tricks}.
W eksperymentach użyto klasyfikatora nadzorowanego z n-gramami znakowymi
długości 3--6, unigramami słów, wektorami o wymiarze 100, współczynnikiem
uczenia 0,25 i 20 epokami treningu.

Drugim wariantem pośrednim jest użycie gotowego modelu embeddingowego do
zamiany wiadomości na wektor semantyczny, a następnie wytrenowanie prostego
klasyfikatora liniowego. Sentence-BERT został zaproponowany właśnie po to, aby
uzyskiwać reprezentacje zdań nadające się do porównań semantycznych i dalszych
zadań bez kosztownego porównywania każdej pary tekstów przez pełny model BERT
~\autocite{reimers-sentence-bert}. W eksperymentach sprawdzono ten wariant za
pomocą MiniLM i regresji logistycznej.

Słabszym punktem takiego wariantu jest brak treningu pod konkretny typ nadużyć.
Gotowy model embeddingowy może dobrze oddawać ogólne znaczenie tekstu, ale
słabiej reagować na drobne sygnały istotne dla filtra pocztowego: nietypowe
domeny, zapis kwot, celowe literówki albo powtarzalne elementy szablonów
kampanii.

\section{Modele transformerowe i językowe}

Modele transformerowe budują reprezentację zależną od kontekstu. W wariancie
klasyfikacyjnym model, taki jak BERT lub DistilBERT, otrzymuje tekst wiadomości
i zwraca etykietę przez głowicę klasyfikacyjną. DistilBERT jest mniejszą wersją
BERT-a uzyskaną przez destylację wiedzy; według autorów zachowuje dużą część
jakości modelu bazowego, przy mniejszym rozmiarze i szybszym działaniu
~\autocite{sanh-distilbert}.

Modele językowe wymagają więcej pamięci niż metody TF-IDF, a ich inferencja jest
wolniejsza i zależy od długości kontekstu. Mogą jednak lepiej wykorzystywać
kontekst i uogólniać poza dosłowne słownictwo zbioru treningowego. Dlatego w
części badawczej oceniono również ich profil błędów na danych spoza zbioru
treningowego.

\subsection{Badania nad wykorzystaniem modeli językowych w filtrowaniu poczty}

W badaniach nad użyciem modeli językowych do filtrowania poczty stosowano
dostrajanie na oznaczonych wiadomościach, klasyfikację na podstawie promptu oraz
analizę treści uzupełnionej o nagłówki lub reputację adresów URL.

W pracy \emph{Spam-T5: Benchmarking Large Language Models for Few-Shot Email
Spam Detection} porównano klasyczne klasyfikatory, RoBERTę, SetFit oraz Spam-T5,
czyli Flan-T5 trenowany do generowania etykiety wiadomości. Eksperyment
obejmował cztery korpusy i próby treningowe od 4 przykładów do pełnego zbioru.
Spam-T5 uzyskał najlepsze wyniki przy 4--16 przykładach oraz najwyższy średni
F1 liczony dla wszystkich badanych rozmiarów próby. Trening i inferencja trwały
dłużej niż w przypadku klasycznych metod~\autocite{labonne-spam-t5-2023}.
W pracy \emph{Next-Generation Spam Filtering: Comparative Fine-Tuning of LLMs,
NLPs, and CNN Models for Email Spam Classification} dostrojono GPT-4, BERT-a i
RoBERTę na dwóch zbiorach wiadomości. Po losowym podziale danych modele osiągały
accuracy przekraczające 98\%. Ocena na drugim korpusie przyniosła niższe wyniki,
szczególnie dla BERT-a i RoBERT-y
~\autocite{roumeliotis-next-generation-spam-2024}.

Klasyfikacja przez prompt wykorzystuje gotowy model wraz z instrukcją i
przykładami umieszczonymi w kontekście. W pracy \emph{Evaluating the Performance
of ChatGPT for Spam Email Detection} sprawdzono ChatGPT w trybie zero-shot,
one-shot i five-shot. Na większym zbiorze anglojęzycznym najlepszy wariant
ChatGPT uzyskał macro F1 równe 80\%, a nadzorowany BERT 98\%. Na małym zbiorze
chińskojęzycznym różnica między ChatGPT a metodami nadzorowanymi była mniejsza
~\autocite{si-chatgpt-spam-2024}. W pracy \emph{Zero-Shot Spam Email
Classification Using Pre-trained Large Language Models} oceniono ChatGPT, GPT-4
i Flan-T5 na całym korpusie SpamAssassin. Modele klasyfikowały skróconą treść
albo streszczenie utworzone wcześniej przez ChatGPT. Etap streszczania podniósł
F1 GPT-4 z 88\% do 95\% i wymagał dodatkowego wywołania modelu. Badanie objęło
jeden korpus, SpamAssassin~\autocite{rojas-galeano-zero-shot-spam-2024}.

W pracy \emph{ChatSpamDetector: Leveraging Large Language Models for Effective
Phishing Email Detection} opisano system przetwarzający nagłówki i elementy
HTML, który oprócz klasy zwracał wyjaśnienie decyzji. GPT-4 osiągnął accuracy
równe 99,70\%. Średni czas jednego zapytania wyniósł 12,53~s, a ocena 2010
wiadomości kosztowała 265,80~USD według cen API użytych w badaniu
~\autocite{koide-chatspamdetector-2024}. W pracy \emph{APOLLO: A GPT-Based Tool
to Detect Phishing Emails and Generate Explanations That Warn Users} połączono
GPT-4o z reputacją adresów URL. Poprawne dane zewnętrzne podniosły accuracy z
97,4\% do 99,9\%. Po podaniu błędnych informacji o adresach URL wynik spadł
poniżej 46\%~\autocite{desolda-apollo-2025}.

Wartości zestawione w tabeli~\ref{tab:related-llm-email-detection} pochodzą z
różnych zbiorów, podziałów danych i zakresów informacji wejściowych. Ich
bezpośrednie porównanie ma charakter orientacyjny.

\begin{table}[!htbp]
\centering
\scriptsize
\resizebox{\textwidth}{!}{\begin{tabular}{p{0.18\textwidth}p{0.21\textwidth}p{0.23\textwidth}p{0.31\textwidth}}
\hline
\textbf{Praca} & \textbf{Modele i sposób użycia} & \textbf{Dane} & \textbf{Najważniejszy wynik i ograniczenie} \\
\hline
Labonne i Moran~\autocite{labonne-spam-t5-2023} & Spam-T5, RoBERTa, SetFit i metody klasyczne; trening pełny i few-shot & Ling-Spam, SMS Spam Collection, SpamAssassin i Enron & F1 Spam-T5 95--99\% przy pełnym treningu; przewaga przy 4--16 przykładach, lecz wyższy czas treningu i inferencji \\
Si i in.~\autocite{si-chatgpt-spam-2024} & ChatGPT zero-, one- i five-shot; BERT oraz metody klasyczne & angielski ESD i mały zbiór chińskojęzyczny & macro F1 80\% dla ChatGPT i 98\% dla BERT-a na ESD; odwrotna relacja na małym zbiorze chińskim \\
Rojas-Galeano~\autocite{rojas-galeano-zero-shot-spam-2024} & ChatGPT, GPT-4 i Flan-T5 bez dostrajania & SpamAssassin, 6047 wiadomości & F1 90\% dla Flan-T5 na treści i 95\% dla GPT-4 na streszczeniach; jeden korpus i dodatkowy etap generowania \\
Koide i in.~\autocite{koide-chatspamdetector-2024} & GPT-4, GPT-3.5, Llama~2 i Gemini Pro; treść oraz nagłówki & 1010 wiadomości phishingowych i 1000 pożądanych & accuracy 99,70\% dla GPT-4; średnio 12,53~s na wiadomość i 265,80~USD za ocenę zbioru \\
Desolda i in.~\autocite{desolda-apollo-2025} & GPT-4o, treść wiadomości i dane o adresach URL & 2000 wiadomości phishingowych i 2000 pożądanych & accuracy 97,4\% bez danych zewnętrznych i 99,9\% z poprawną reputacją URL; silna zależność od jakości tych danych \\
Roumeliotis i in.~\autocite{roumeliotis-next-generation-spam-2024} & dostrajanie GPT-4, BERT-a, RoBERT-y i CNN & dwa zbiory z serwisu Kaggle & accuracy powyżej 98\% na własnych testach; spadek przy ocenie na drugim zbiorze \\
\hline
\end{tabular}
}
\caption{Badania nad wykorzystaniem modeli transformerowych i językowych do
klasyfikacji spamu oraz phishingu w poczcie elektronicznej.}
\label{tab:related-llm-email-detection}
\end{table}
\FloatBarrier

W analizowanych pracach małe modele generatywne z otwartymi wagami, liczące
poniżej miliarda parametrów, pojawiają się rzadko. Badania skupiają się zwykle
na jednym sposobie użycia modelu. Spośród zestawionych prac tylko
\emph{Next-Generation Spam Filtering} obejmuje test przenoszenia modelu między
korpusami.

Wkład własny pracy obejmuje porównanie czterech sposobów realizacji
klasyfikacji przez Qwen3-0.6B przy zachowaniu wspólnego modelu, danych i
protokołu oceny. Porównanie przeprowadzono na kilku korpusach, a wyniki
zestawiono z prostszymi klasyfikatorami oraz kosztem inferencji wszystkich
metod.

\section{Zakres porównania eksperymentalnego}

Bezpośrednie porównanie objęło metody korzystające z treści wiadomości.
Wiadomości przypisywano do klasy \texttt{ham} albo \texttt{spam}. Wykorzystane
zbiory nie zawierają spójnych danych o reputacji nadawcy, DNS, pełnych
nagłówkach ani historii ruchu.
Tabela~\ref{tab:sota-method-scope} wskazuje, które metody omówiono jako kontekst
literaturowy, a które uruchomiono w eksperymentach.

\begin{table}[H]
\centering
\small
\resizebox{\textwidth}{!}{\begin{tabular}{p{0.22\textwidth}p{0.31\textwidth}p{0.18\textwidth}p{0.18\textwidth}}
\hline
\textbf{Grupa metod} & \textbf{Główne źródło sygnału} & \textbf{Wymaga metadanych poczty} & \textbf{Ujęcie w eksperymentach} \\
\hline
Uwierzytelnianie i reputacja nadawcy & Nagłówki SMTP, DNS, SPF, DKIM, DMARC, reputacja adresów IP i domen & Tak & Nie \\
Systemy regułowe i reputacyjne & Reguły treści, adresy URL, listy reputacyjne, scoring, czasem filtr bayesowski & Częściowo & Nie \\
TF-IDF + klasyfikator & Treść wiadomości reprezentowana przez cechy leksykalne & Nie & Tak \\
\texttt{fastText} & N-gramy słów i szybki model liniowy z embeddingami & Nie & Tak \\
Embeddingi zdań + klasyfikator & Wektor semantyczny wiadomości i klasyfikator liniowy & Nie & Tak \\
Transformer klasyfikacyjny & Dostrajany enkoder tekstu z głowicą klasyfikacyjną & Nie & Tak \\
Mały model językowy z LoRA & Dostrojony model generatywny, etykieta lub prawdopodobieństwo etykiety & Nie & Główna metoda opisywana w pracy \\
\hline
\end{tabular}
}
\caption{Zakres metod omawianych w przeglądzie oraz ich wykorzystanie w części
eksperymentalnej.}
\label{tab:sota-method-scope}
\end{table}
